% TOP_PROBLEM: {"id": 1428, "title": "Vizing's Conjecture", "category_id": 3, "category": {"id": 3, "name": "graph_theory", "display_name": "Graph Theory", "description": "Problems involving graphs, networks, and their properties.", "slug": "graph-theory", "order_index": 3, "created_at": "2026-07-31T15:26:25.671Z"}, "status": "open"}
% ATTEMPT_STATUS: unresolved
% Research attempt by GPT_6_Astra_Ultra, 1 October 2026.
% Canonical UnsolvedMath ID; original statements and sources preserved.
% FIRST_POSTED: 2026-10-01
% Imported from: unsolvedmath_additional_500.tex; UnsolvedMath ID 1428
% !TeX program = lualatex
% Compile twice: lualatex 1428.tex
% Self-contained: no external bibliography, images, or \input files.
% Numeric section labels and visible numbers are original UnsolvedMath IDs.
\documentclass[11pt,a4paper]{article}
\usepackage[margin=24mm,headheight=14pt]{geometry}
\usepackage{fontspec}
\setmainfont{Latin Modern Roman}
\setsansfont{Latin Modern Sans}
\setmonofont{Latin Modern Mono}
\usepackage{amsmath,amssymb,mathtools}
\usepackage{unicode-math}
\setmathfont{Latin Modern Math}
\usepackage{microtype}
\usepackage{enumitem,xurl,needspace}
\usepackage[dvipsnames]{xcolor}
\usepackage{hyperref}
\hypersetup{unicode=true,colorlinks=true,linkcolor=MidnightBlue,urlcolor=MidnightBlue,
 pdftitle={UnsolvedMath Problem 1428},pdfauthor={GPT 6 Astra Ultra},bookmarksnumbered=true}
\usepackage{bookmark,tocloft,titlesec,fancyhdr}
\setlength{\cftsecnumwidth}{4.2em}
\cftsetpnumwidth{2.5em}
\cftsetrmarg{4em}
\titleformat{\section}{\Large\bfseries\raggedright}{\thesection}{0.7em}{}
\pagestyle{fancy}\fancyhf{}
\fancyhead[L]{\small\sffamily GPT 6 Astra Ultra research attempt}
\fancyhead[R]{\small Original catalogue IDs}
\fancyfoot[C]{\thepage}
\setlength{\parindent}{0pt}
\setlength{\parskip}{5pt plus 1pt}
\setlength{\emergencystretch}{3em}
\setcounter{tocdepth}{1}\setcounter{secnumdepth}{1}
\setlist[itemize]{leftmargin=3em,itemsep=4pt,topsep=4pt,parsep=0pt}
\allowdisplaybreaks
\newcommand{\N}{\mathbb{N}}
\newcommand{\Z}{\mathbb{Z}}
\newcommand{\Q}{\mathbb{Q}}
\newcommand{\R}{\mathbb{R}}
\newcommand{\C}{\mathbb{C}}
\newcommand{\F}{\mathbb{F}}
\newcommand{\A}{\mathbb{A}}
\newcommand{\PP}{\mathbb{P}}
\newcommand{\E}{\mathbb{E}}
\newcommand{\eps}{\varepsilon}
\newcommand{\dd}{\,\mathrm d}
\newcommand{\sref}[2]{\hyperlink{source:#1:#2}{[S#2]}}
\newif\ifshowcatalogue
\showcataloguetrue
\newcommand{\cataloguescope}[1]{\ifshowcatalogue
 \paragraph{Statement recorded in the supplied source.}
 #1\fi}
\begin{document}
% UnsolvedMath ID 1428; source catalogue code GRAPH-041

\setcounter{section}{1427}

\section{Vizing’s Cartesian-Product Domination Conjecture}\label{1428}

{\small\sffamily UnsolvedMath ID 1428 · Graph Theory}

\subsection{Definitions and mathematical statement}

\paragraph{Shared notation.}Unless a section says otherwise, $\N=\{1,2,\ldots\}$,
$\Z$ is the ring of integers, $\Q,\R,\C$ are the rational, real, and complex
fields, $[N]=\{1,\ldots,\lfloor N\rfloor\}$, and $\log$ is the natural
logarithm. For real functions of a parameter tending to infinity,
$f\ll g$ and $f=O(g)$ mean $|f|\leq Cg$ eventually for some fixed $C$;
$f\gg g$ reverses this comparison; $f=o(g)$ means $f/g\to0$;
$f\sim g$ means $f/g\to1$; and $f\asymp g$ means both order bounds.
A subscript on $O$ or $\ll$ specifies allowed constant dependence.
The expression $n^{o(1)}$ in an upper bound means growth smaller than
$n^\epsilon$ for each fixed $\epsilon>0$. Local definitions govern any
exception, finite-field convention, or use of the same symbol for another
object. An asymptotic claim is not automatically an effective algorithm.



A dominating set meets the closed neighbourhood of every vertex, and $\gamma(G)$ is its minimum size. In $G\square H$, vertices are pairs $(g,h)$, with adjacency when one coordinate agrees and the other is an edge of the corresponding factor. The conjecture is $\gamma(G\square H)\geq\gamma(G)\gamma(H)$ for finite graphs. \sref{1428}{1} \sref{1428}{2}

\paragraph{Statement recorded in the supplied source.}
For the Cartesian product of graphs $G \square H$, is the domination number at least $\gamma(G) \cdot \gamma(H)$? \sref{1428}{1}

\subsection{Short English statement}

Must dominating a Cartesian product require at least the product of the minimum dominating-set sizes of its two factors?

\subsection{Research attempt}

\paragraph{Target and selected packing route.}
The domination number uses closed neighbourhoods and the product is
Cartesian. The recent primary records [E1,E2] discuss constant-factor
progress; their titles and bounds do not give the full coefficient one
for all graphs. This attempt proves a packing bound that implies the
conjecture for a substantial subclass, then isolates the overlap that
prevents the same count from handling arbitrary factors.

\paragraph{Projection supplies only the individual lower bounds.}
Let $D$ dominate $G\square H$. Its first-coordinate projection
dominates $G$: otherwise there is a vertex $g$ with no projected
coordinate in $N_G[g]$, and no $(g,h)$ could be dominated in the
product. Thus $|D|\ge\gamma(G)$, and symmetrically
\[
                 \gamma(G\square H)\ge\max\{\gamma(G),\gamma(H)\}.
\]
The two projections come from the same elements of $D$. Multiplying
their separate lower bounds is invalid; a set of ordered pairs may
have large projections without containing their Cartesian product.
This identifies why the simplest projection argument stops short.

\paragraph{Disjoint closed neighbourhoods give a stronger theorem.}
Let $g_1,\ldots,g_r$ have pairwise disjoint closed neighbourhoods
in $G$. Such a set is often called a distance-two packing, because
its distinct centres have distance at least three whenever connected.
For each $i$ define
\[
            D_i=D\cap(N_G[g_i]\times V(H)).
\]
The second-coordinate projection of $D_i$ dominates $H$. Indeed,
to dominate $(g_i,h)$ a member of $D$ must either have first
coordinate $g_i$ and second coordinate in $N_H[h]$, or have first
coordinate adjacent to $g_i$ and second coordinate exactly $h$.
In either case it belongs to $D_i$ and its second coordinate lies
in $N_H[h]$. Hence $|D_i|\ge\gamma(H)$.

The neighbourhoods are disjoint, so the $D_i$ are disjoint subsets
of $D$. Summing gives
\[
                     |D|\ge r\gamma(H).
\]
If $\rho(G)$ denotes the largest size of such a packing, we obtain
\[
              \gamma(G\square H)\ge\rho(G)\gamma(H),
              \qquad
              \gamma(G\square H)\ge\gamma(G)\rho(H).
\]
Also $\rho(G)\le\gamma(G)$: a dominating set must meet every
$N_G[g_i]$, and those sets are disjoint. Thus the bound proves the
desired conjecture whenever one factor satisfies $\rho=\gamma$.

\paragraph{Efficient domination gives a complete subclass.}
Suppose some centres $g_1,\ldots,g_r$ have closed neighbourhoods
partitioning $V(G)$. They form a dominating set, so $\gamma(G)\le r$;
they also form a packing, so $\gamma(G)\ge\rho(G)\ge r$.
Consequently $\rho(G)=\gamma(G)=r$, and the preceding theorem
establishes Vizing's inequality for this $G$ and every finite graph
$H$. This argument requires no additional properties of $H$.

For the path $P_n$ with vertices $1,\ldots,n$, explicit such centres
are as follows. If $n=3k$, use $2,5,\ldots,3k-1$. If $n=3k+1$,
use $1,4,\ldots,3k+1$. If $n=3k+2$, use $2,5,\ldots,3k+2$.
The end neighbourhoods are truncated at the endpoints and the interior
ones are consecutive three-vertex intervals. They are disjoint and
cover the path in each case, including $n=1,2$. Therefore
\[
                   \gamma(P_n)=\lceil n/3\rceil,
 \qquad \gamma(P_n\square H)\ge\lceil n/3\rceil\gamma(H).
\]
Likewise in $C_{3k}$, $k\ge1$, centres $1,4,\ldots,3k-2$
have disjoint three-vertex closed neighbourhoods covering the cycle.
Thus $\gamma(C_{3k})=k$ and the product conjecture holds with such
a cycle as one factor. For a disjoint union of efficiently dominated
components, taking all their centres preserves the same property.

\paragraph{A precise obstruction to the unrestricted counting argument.}
For $C_4$, the domination number is two: no closed neighbourhood has
all four vertices, while two opposite vertices dominate. However any
two vertices of $C_4$ are at distance at most two, so $\rho(C_4)=1$.
Thus the packing bound gives only $\gamma(H)$ when the desired bound
is $2\gamma(H)$. The gap is already present in a four-vertex factor.

Starting with a minimum dominating set rather than a packing does not
fix the sum. Its closed neighbourhoods may overlap, causing a single
element of $D$ to be counted in several sets $D_i$. Dividing by a
maximum overlap multiplicity gives a weaker bound, and no argument
here makes that multiplicity one. To reach the conjecture for general
factors one needs a more careful allocation or charging scheme that
preserves a full $\gamma(H)$ contribution per dominating centre.
The elementary packing proof supplies no such allocation.

\paragraph{Executed finite tests.}
The script enumerates vertex subsets in increasing cardinality until
a dominating set is found, using exact closed-neighbourhood bit masks.
Testing every smaller subset certifies minimality for each graph.
For $P_2\square P_3$, $P_4\square P_4$, and $C_4\square C_4$
it finds domination numbers $2,4,4$, respectively. The corresponding
products of the factor domination numbers are $1,4,4$. Thus the
inequality can be strict and can also attain equality; the $C_4$
example succeeds despite the weakness of the packing bound.
Witness sets are saved with the vertex encoding used by the script.

These computations are in \texttt{checks\_graphs.py} and
\texttt{results\_graphs.json} under
\path{_Local/Erdos_problems/UnsolvedMath/attempt_audit_2026_10_01/number_theory/}.
They are checks of three finite products, not an enumeration of all
finite graphs or a replication of the recent constant-factor papers.

\paragraph{Remaining gap and outcome.}
The packing theorem proves the full inequality whenever one factor
has equal packing and domination numbers, including all paths and
cycles of order divisible by three. Overlapping neighbourhoods in a
general minimum dominating set prevent the same proof from extending.
No general allocation theorem or counterexample is obtained.
\textbf{Outcome: UNRESOLVED.}

\subsection{Sources}

{\small

\begin{itemize}

\item[\hypertarget{source:1428:1}{S1}] ULAM.AI, \emph{UnsolvedMath}, supplied \texttt{problems.json}, record 1428 (GRAPH-041), statement and background. The embedded literature-review date is 2026-08-17; that is the archive’s review date, not a fresh certification. Dataset landing page: \url{https://huggingface.co/datasets/ulamai/UnsolvedMath}.

\item[\hypertarget{source:1428:2}{S2}] Stephen Suen and Jennifer Tarr, An improved inequality related to Vizing's conjecture, Electronic Journal of Combinatorics 19 (2012), \#P8. \url{https://doi.org/10.37236/15}. Bibliographic information imported from the supplied record.

\item[\hypertarget{source:1428:3}{S3}] Raphael Steiner, A constant-factor step towards Vizing's conjecture, arXiv:2606.14414 (2026). \url{https://arxiv.org/abs/2606.14414}. Bibliographic information imported from the supplied record.

\item[\hypertarget{source:1428:4}{S4}] Mohsen Aliabadi and Elliot Krop, An improved constant for Vizing's conjecture, arXiv:2607.01109 (2026). \url{https://arxiv.org/abs/2607.01109}. Bibliographic information imported from the supplied record.


\item[E1] Raphael Steiner, \emph{A constant-factor
step towards Vizing's conjecture}, arXiv:2606.14414.
\url{https://arxiv.org/abs/2606.14414}. Primary record inspected
1 October 2026; background only, not a premise of the packing proof.
\item[E2] Mohsen Aliabadi and Elliot Krop, \emph{A $2/3$ Bound for
Vizing's Conjecture}, arXiv:2607.01109.
\url{https://arxiv.org/abs/2607.01109}. Primary title and record
inspected 1 October 2026. This current title differs from the imported
bibliographic wording in [S4]; a coefficient below one is not the
full inequality recorded in this problem.

\end{itemize}

}

\label{end:1428}
\end{document}
